% GENERATED FILE. Edit canonical lessons, then run npm run generate:books.
% canonical-source-hash: fnv1a64:21fced49d6f28952
% canonical-lessons: TA-C465-enku, TA-C465-pacam, TA-C465-piriyam, TA-C465-mokam, TA-C465-nekil

\chapter{Longing and Fondness}
\label{ch:ta-longing-and-fondness}
\hlchaptermodality{465}

\begin{chapteropening}
\textbf{By the end of this chapter:} \emph{I can talk about longing, family affection, fondness and being moved.}

Complete the last lesson of chapter 465: I can talk about longing, family affection, fondness and being moved.
\end{chapteropening}

\section[\texorpdfstring{\ta{ஏங்கு} (ēṅku)}{ஏங்கு (ēṅku)}]{\ta{ஏங்கு} (ēṅku) — to long for; to pine}

\label{lesson:TA-C465-enku}



\begin{quote}
Before the new one: say the Tamil for to plead, then the Tamil for to be fed up.
\end{quote}

\subsection*{You'll want to know: \ta{ஏங்கு}}
\textbf{\ta{ஏங்கு}} — \emph{ēṅku} — \textquotedblleft{}to long for; to pine\textquotedblright{}.

\begin{quote}
\ta{வெளிநாட்டில்} \ta{இருந்தபோது} \ta{நான்} \ta{அம்மாவின்} \ta{சமையலுக்கு} \ta{ஏங்கினேன்}.
\end{quote}

\begin{quote}
\emph{veḷināṭṭil iruntapōtu nāṉ ammāviṉ camaiyalukku ēṅkiṉēṉ.}
\end{quote}

\begin{quote}
When I was abroad I longed for my mother’s cooking.
\end{quote}

Past (I …): \ta{ஏங்கினேன்}, \emph{ēṅkiṉēṉ}.

What one longs for takes -\ta{க்கு}: \ta{அன்புக்கு} \ta{ஏங்கு}, to crave love.

\begin{cousinweb}
\textbf{\ta{ஏங்கு}} (\emph{ēṅku}) is native Dravidian (DEDR 878); DEDR keeps a second \ta{ஏங்கு}, \textquotedblleft{}to cry out\textquotedblright{}, as its own entry (879).
\end{cousinweb}

\begin{grammarlens}[title={What you've built}]
One more word for describing what you lived and felt.
\end{grammarlens}

\subsection*{Your turn}
\begin{itemize}
\raggedright
  \item \emph{Say it:} \emph{ēṅku}
  \item \emph{Say it:} \emph{ēṅku}, once more
  \item \emph{Say it:} \emph{veḷināṭṭil iruntapōtu nāṉ ammāviṉ camaiyalukku ēṅkiṉēṉ.} — When I was abroad I longed for my mother’s cooking.
  \item \emph{Recall:} say the Tamil for to plead, then the Tamil for to be fed up, then say \emph{ēṅku} again
\end{itemize}

\begin{tcolorbox}[breakable,colback=teal!4,colframe=teal!35!black,title={Before you move on}]
What does \ta{ஏங்கு} mean? (\textquotedblleft{}To long for; to pine\textquotedblright{}.) Say it once more.
\end{tcolorbox}
\section[\texorpdfstring{\ta{பாசம்} (pācam)}{பாசம் (pācam)}]{\ta{பாசம்} (pācam) — family affection; attachment}

\label{lesson:TA-C465-pacam}



\begin{quote}
Before the new one: say the Tamil for to be fed up, then the Tamil for to long for.
\end{quote}

\subsection*{You'll want to know: \ta{பாசம்}}
\textbf{\ta{பாசம்}} — \emph{pācam} — \textquotedblleft{}family affection; attachment\textquotedblright{}.

\begin{quote}
\ta{தாத்தாவுக்குப்} \ta{பேரன்} \ta{மேல்} \ta{நிறைய} \ta{பாசம்}.
\end{quote}

\begin{quote}
\emph{tāttāvukkup pēraṉ mēl niṟaiya pācam.}
\end{quote}

\begin{quote}
Grandfather is very attached to his grandson.
\end{quote}

\ta{மேல்} \textquotedblleft{}on\textquotedblright{} marks the one loved: \ta{அவன்} \ta{மேல்} \ta{பாசம்}.

\begin{cousinweb}
\textbf{\ta{பாசம்}} (\emph{pācam}) is borrowed from Sanskrit \emph{pāśa}, \textquotedblleft{}a noose, a cord\textquotedblright{} (CDIAL 8133): in Tamil, the tie of love that binds a family.
\end{cousinweb}

\begin{grammarlens}[title={What you've built}]
One more word for describing what you lived and felt.
\end{grammarlens}

\subsection*{Your turn}
\begin{itemize}
\raggedright
  \item \emph{Say it:} \emph{pācam}
  \item \emph{Say it:} \emph{pācam}, once more
  \item \emph{Say it:} \emph{tāttāvukkup pēraṉ mēl niṟaiya pācam.} — Grandfather is very attached to his grandson.
  \item \emph{Recall:} say the Tamil for to be fed up, then the Tamil for to long for, then say \emph{pācam} again
\end{itemize}

\begin{tcolorbox}[breakable,colback=teal!4,colframe=teal!35!black,title={Before you move on}]
What does \ta{பாசம்} mean? (\textquotedblleft{}Family affection; attachment\textquotedblright{}.) Say it once more.
\end{tcolorbox}
\section[\texorpdfstring{\ta{பிரியம்} (piriyam)}{பிரியம் (piriyam)}]{\ta{பிரியம்} (piriyam) — fondness; liking}

\label{lesson:TA-C465-piriyam}



\begin{quote}
Before the new one: say the Tamil for to long for, then the Tamil for family affection.
\end{quote}

\subsection*{You'll want to know: \ta{பிரியம்}}
\textbf{\ta{பிரியம்}} — \emph{piriyam} — \textquotedblleft{}fondness; liking\textquotedblright{}.

\begin{quote}
\ta{என்} \ta{தங்கைக்கு} \ta{மல்லிகைப்} \ta{பூ} \ta{என்றால்} \ta{பிரியம்}.
\end{quote}

\begin{quote}
\emph{eṉ taṅkaikku mallikaip pū eṉṟāl piriyam.}
\end{quote}

\begin{quote}
My younger sister is fond of jasmine.
\end{quote}

X \ta{என்றால்} \ta{பிரியம்}, to be fond of X.

\begin{cousinweb}
\textbf{\ta{பிரியம்}} (\emph{piriyam}) is borrowed from Sanskrit \emph{priya}, \textquotedblleft{}dear, beloved\textquotedblright{} (CDIAL 8974); Tamil breaks the \emph{pr} with a vowel: \ta{பிரி}-.
\end{cousinweb}

\begin{grammarlens}[title={What you've built}]
One more word for describing what you lived and felt.
\end{grammarlens}

\subsection*{Your turn}
\begin{itemize}
\raggedright
  \item \emph{Say it:} \emph{piriyam}
  \item \emph{Say it:} \emph{piriyam}, once more
  \item \emph{Say it:} \emph{eṉ taṅkaikku mallikaip pū eṉṟāl piriyam.} — My younger sister is fond of jasmine.
  \item \emph{Recall:} say the Tamil for to long for, then the Tamil for family affection, then say \emph{piriyam} again
\end{itemize}

\begin{tcolorbox}[breakable,colback=teal!4,colframe=teal!35!black,title={Before you move on}]
What does \ta{பிரியம்} mean? (\textquotedblleft{}Fondness; liking\textquotedblright{}.) Say it once more.
\end{tcolorbox}
\section[\texorpdfstring{\ta{மோகம்} (mōkam)}{மோகம் (mōkam)}]{\ta{மோகம்} (mōkam) — infatuation; craze}

\label{lesson:TA-C465-mokam}



\begin{quote}
Before the new one: say the Tamil for family affection, then the Tamil for fondness.
\end{quote}

\subsection*{You'll want to know: \ta{மோகம்}}
\textbf{\ta{மோகம்}} — \emph{mōkam} — \textquotedblleft{}infatuation; craze\textquotedblright{}.

\begin{quote}
\ta{என்} \ta{மகனுக்குக்} \ta{கிரிக்கெட்} \ta{மேல்} \ta{மோகம்}.
\end{quote}

\begin{quote}
\emph{eṉ makaṉukkuk kirikkeṭ mēl mōkam.}
\end{quote}

\begin{quote}
My son is crazy about cricket.
\end{quote}

A passion that clouds judgment: \ta{பண} \ta{மோகம்}, a craving for money.

\begin{cousinweb}
\textbf{\ta{மோகம்}} (\emph{mōkam}) is borrowed from Sanskrit \emph{mōha}, \textquotedblleft{}bewilderment; a swoon\textquotedblright{} (CDIAL 10360).
\end{cousinweb}

\begin{grammarlens}[title={What you've built}]
One more word for describing what you lived and felt.
\end{grammarlens}

\subsection*{Your turn}
\begin{itemize}
\raggedright
  \item \emph{Say it:} \emph{mōkam}
  \item \emph{Say it:} \emph{mōkam}, once more
  \item \emph{Say it:} \emph{eṉ makaṉukkuk kirikkeṭ mēl mōkam.} — My son is crazy about cricket.
  \item \emph{Recall:} say the Tamil for family affection, then the Tamil for fondness, then say \emph{mōkam} again
\end{itemize}

\begin{tcolorbox}[breakable,colback=teal!4,colframe=teal!35!black,title={Before you move on}]
What does \ta{மோகம்} mean? (\textquotedblleft{}Infatuation; craze\textquotedblright{}.) Say it once more.
\end{tcolorbox}
\section[\texorpdfstring{\ta{நெகிழ்} (nekiḻ)}{நெகிழ் (nekiḻ)}]{\ta{நெகிழ்} (nekiḻ) — to be moved; to be touched}

\label{lesson:TA-C465-nekil}



\begin{quote}
Before the new one: say the Tamil for fondness, then the Tamil for infatuation.
\end{quote}

\subsection*{You'll want to know: \ta{நெகிழ்}}
\textbf{\ta{நெகிழ்}} — \emph{nekiḻ} — \textquotedblleft{}to be moved; to be touched\textquotedblright{}.

\begin{quote}
\ta{அவர்களுடைய} \ta{அன்பைப்} \ta{பார்த்து} \ta{நான்} \ta{நெகிழ்ந்தேன்}.
\end{quote}

\begin{quote}
\emph{avarkaḷuṭaiya aṉpaip pārttu nāṉ nekiḻntēṉ.}
\end{quote}

\begin{quote}
I was moved by their affection.
\end{quote}

Past (I …): \ta{நெகிழ்ந்தேன்}, \emph{nekiḻntēṉ}.

\ta{மனம்} \ta{நெகிழ்ந்தது}, my heart was touched.

\begin{cousinweb}
\textbf{\ta{நெகிழ்}} (\emph{nekiḻ}) is native Dravidian (DEDR 2923); the entry runs from loosening, as bangles slip, to a softened heart.
\end{cousinweb}

\begin{grammarlens}[title={What you've built}]
One more word for describing what you lived and felt.
\end{grammarlens}

\subsection*{Your turn}
\begin{itemize}
\raggedright
  \item \emph{Say it:} \emph{nekiḻ}
  \item \emph{Say it:} \emph{nekiḻ}, once more
  \item \emph{Say it:} \emph{avarkaḷuṭaiya aṉpaip pārttu nāṉ nekiḻntēṉ.} — I was moved by their affection.
  \item \emph{Recall:} say the Tamil for fondness, then the Tamil for infatuation, then say \emph{nekiḻ} again
\end{itemize}

\begin{tcolorbox}[breakable,colback=teal!4,colframe=teal!35!black,title={Before you move on}]
What does \ta{நெகிழ்} mean? (\textquotedblleft{}To be moved; to be touched\textquotedblright{}.) Say it once more.
\end{tcolorbox}
